% ============================================================
%  TOM TAT DO AN
% ============================================================
\unnumberedchapter{TÓM TẮT ĐỒ ÁN}

Việc quyết định ``hôm nay ăn gì'' là công việc lặp lại hằng ngày của mỗi gia
đình: người nội trợ phải tự ghi nhớ nguyên liệu đang có, tự nghĩ món phù hợp,
rồi tự lập danh sách đi chợ cho cả tuần. Ba việc này hiện chưa được hỗ trợ một
cách liền mạch bởi các công cụ phổ biến tại Việt Nam. Đồ án xây dựng
\textbf{CookingAdvisor}, một ứng dụng web giải quyết đồng thời cả ba việc nêu
trên, phát triển trên nền tảng \textbf{ASP.NET Core MVC (.NET 10)}, sử dụng
\textbf{Entity Framework Core} theo hướng Code-First với hệ quản trị cơ sở dữ
liệu \textbf{SQL Server}, và \textbf{ASP.NET Core Identity} cho xác thực, phân
quyền theo hai vai trò Quản trị viên và Người dùng.

Đóng góp kỹ thuật của đồ án nằm ở ba thuật toán trong tầng nghiệp vụ:
\textbf{thuật toán gợi ý theo nguyên liệu sẵn có} tính độ phủ giữa tập nguyên
liệu người dùng đang có và tập nguyên liệu của từng món, xếp hạng theo thứ tự
ưu tiên món nấu được ngay, độ phủ giảm dần, số nguyên liệu còn thiếu ít nhất;
\textbf{thuật toán sinh thực đơn tuần} theo hướng tham lam có ràng buộc, lấp đầy
21 suất ăn (7 ngày $\times$ 3 bữa) với mục tiêu tránh lặp món, tôn trọng ràng
buộc vùng miền và cân đối năng lượng theo ngày; \textbf{thuật toán sinh danh
sách đi chợ} gộp nguyên liệu của thực đơn theo cặp (nguyên liệu, đơn vị) và cộng
dồn khối lượng.

Kết quả là một hệ thống hoàn chỉnh chạy đầu-cuối với cơ sở dữ liệu mẫu gồm 26
món ăn Việt Nam, 40 nguyên liệu, 10 danh mục, cung cấp 20 chức năng phân theo ba
mức quyền truy cập. Bộ \textbf{29 ca kiểm thử đơn vị} viết bằng xUnit cho ba
thuật toán lõi toàn bộ đạt, đồng thời phát hiện và khắc phục một khiếm khuyết
thực tế: khi thư viện món không đủ lớn, tiêu chí phụ theo năng lượng khiến hệ
thống lặp lại đúng một món cho mọi suất còn trống; sau khi bổ sung tiêu chí rải
đều theo số lần đã sử dụng, thực đơn đạt 17 món khác nhau trên 21 suất. Giao
diện đáp ứng trên dải màn hình từ 375 px đến 1440 px và đạt chuẩn tương phản
WCAG 2.1 mức AA.

\vspace{6pt}
\noindent\textbf{Từ khóa:} gợi ý món ăn, lập thực đơn, độ phủ tập hợp, thuật
toán tham lam, ASP.NET Core MVC, Entity Framework Core.
